\section{Introduction}

Bike-sharing systems depend on local and time-dependent demand. If too few 
bikes are available at a station, users cannot start their trip. If too many 
bikes accumulate, redistribution costs increase and docking space is wasted. 
Accurate demand forecasts can help operators plan bike redistribution in 
advance and improve the overall service quality.

This project explores whether machine learning can produce reliable daily 
demand forecasts at the station level. We use real trip data from a public 
bike-sharing system, combined with weather and calendar information. 
The problem is well-suited for regression because the output is a 
numeric count and the input features are measurable in advance.

Rather than focusing only on the final model score, we treat the full 
pipeline as part of the contribution. We show that preprocessing decisions, 
especially how past demand is represented, have a stronger impact on 
model quality than the choice of algorithm alone.

This report is structured as follows: Section~\ref{sec:dataset} describes 
the dataset and preprocessing steps. Section~\ref{sec:pipeline} explains 
the modelling pipeline and evaluation strategy. Section~\ref{sec:results} 
presents and compares the results of all models. There is also a comparison with another bike-rental project on the same database. 
The remaining sections 
cover ethical considerations, limitations, and a critical reflection.